\section{Methodology}
\label{sec:method}

% Pipeline overview, included from chapters/03_methodology.tex
\begin{figure}[tb]
\centering
\resizebox{\textwidth}{!}{%
\begin{tikzpicture}[
  font=\footnotesize, >=Stealth, node distance=4mm and 5mm,
  box/.style={draw, rounded corners=2pt, align=center, minimum height=9mm, inner sep=3pt, fill=#1},
  box/.default=white]
  \node[box=gray!12] (img) {Field image\\$\approx 4000{\times}3000$};
  \node[box, right=of img] (frame) {Frame detection\\+ perspective crop};
  \node[box, right=of frame] (veg) {Vegetation mask\\$\to$ plant regions};

  % MIL branch (main)
  \node[box=blue!8, right=8mm of veg, yshift=10mm] (dino) {Frozen DINOv3\\per plant $\to h_i$};
  \node[box=blue!8, right=of dino] (head) {Head $f$\\$\hat y_i = f(h_i)$};
  \node[box=blue!15, right=of head] (pool) {Area-weighted mean\\$\hat y = \sum_i w_i\, \hat y_i$};

  % classical branch
  \node[box=orange!10, right=8mm of veg, yshift=-10mm] (direct)
    {Holes / pitting inside leaf mask:\ \ $d = 100\cdot\frac{\text{damaged px}}{\text{leaf px}}$};

  % whole-image baseline
  \node[box=gray!6, below=12mm of img.south west, anchor=north west] (whole)
    {Whole-image baseline: full photo resized to $224{\times}224$ $\to$ frozen DINOv3 $\to$ head $\to \hat y$};

  \draw[->] (img) -- (frame);
  \draw[->] (frame) -- (veg);
  \draw[->] (veg.east) -- ++(4mm,0) |- (dino.west);
  \draw[->] (veg.east) -- ++(4mm,0) |- (direct.west);
  \draw[->] (dino) -- (head);
  \draw[->] (head) -- (pool);
  \draw[->] (img.south) -- (img.south |- whole.north);
  \draw[->, dashed] (veg.north) |- node[pos=0.75, above]{mask areas $A_i$, \ $w_i = A_i / \sum_j A_j$} ([yshift=6mm]pool.north) -- (pool.north);
\end{tikzpicture}}
\caption{Pipeline overview. Blue: area-weighted MIL (main method), trained with the
image-level expert score only. Orange: classical direct damage measurement. Grey: whole-image
baseline.}
\label{fig:pipeline}
\end{figure}


Figure~\ref{fig:pipeline} summarizes the pipeline. We either (i) regress a score from the whole
downsampled photograph (baseline), or reduce the image to the interior of the metal frame, find the
plant regions, and then (ii) measure damage directly with classical colour rules or (iii) regress a
score for each plant region and pool the plant scores into an image score (MIL). All learned
models share the same frozen DINOv3 backbone and are trained on image-level expert scores only.

\subsection{Data curation and frame extraction}
\label{sec:data}

\paragraph{Calibration set.} We use the 470 BBCH10 images from Gro\ss-Gerau (21.10.2025) for
which the independent JLU and GAU scores differ by less than 5 percentage points. The target is
the damaged share of visible leaf area in percent. The images are split by \emph{plot group} into
331 training, 66 validation and 73 test images, so no plot and no physical image appears in two
splits. The split is stored as a fixed manifest, and every feature cache and checkpoint records the
manifest's SHA-256 hash; training and evaluation refuse to run on a mismatching manifest. An
earlier experiment on a different manifest contained duplicated images across splits; its results
are discarded. The test split stayed closed until model selection was finished.

\paragraph{Frame extraction.} Only plants inside the frame are scored. The galvanized frame bars
are segmented by an HSV colour range on a downscaled image, their outer edges are found with a
probabilistic Hough transform, and the four frame corners are estimated from the clustered lines.
The polygon is inset by 2.5\,\% to remove the bars themselves and rectified with a perspective
warp. Images are read with their EXIF orientation applied, since most field images carry a
non-default orientation tag.

\subsection{Classical computer vision}
\label{sec:classical}

\paragraph{Plant regions and area weights.} Inside the frame, a conservative vegetation mask keeps
pixels with OpenCV hue in $[35,85]$, saturation and value $\geq 40$, and excess
green~\cite{woebbecke1995color} $G-\max(R,B) \geq 5$, followed by morphological opening/closing and
removal of components below 30\,px. Nearby green fragments are merged by a dilation proportional to
the image size (1.8\,\% of the shorter side), so that the cotyledons of one seedling form one
region. Regions with fewer than 150 green pixels are dropped and the remaining boxes are padded by
2.5\,\%. If no frame is detected, the whole image is used instead of the frame crop. The number of green
mask pixels $A_i$ of region $i$ is its visible area, used as the MIL weight in
Section~\ref{sec:mil}.

\paragraph{Direct damage measurement.} Following the supervisor's suggestion, we also measured
damage directly as
\begin{equation}
  d = 100 \cdot \frac{\#\text{pixels classified as hole or pitting}}{\#\text{estimated leaf pixels}} .
\end{equation}
For each plant region, the leaf mask is filled to a solid leaf outline, and non-green blobs
enclosed by it are treated as damage: dark blobs as shot holes and bright yellow/brown blobs as
pitting. A second hole rule additionally accepts enclosed blobs whose CIELAB colour is close to
the surrounding soil, since a real hole shows the soil below. Both rules
were judged by manual review (Section~\ref{sec:res-classical}).

\paragraph{Learned hole/pitting detector.} As the colour rules could not separate holes from
pitting and soil, we annotated bounding boxes for \texttt{shot\_hole} and \texttt{pitting} on 40
training/validation images (342 boxes, SAM-assisted in AnyLabeling) and fine-tuned RF-DETR
Nano~\cite{rfdetr}. Because a typical mark is about $18\times18$\,px in a $4000\times3000$ image, we
train on $640\times640$ tiles around the annotations (71 training / 16 validation tiles). This
detector is exploratory and not used by the scoring models.

\subsection{Whole-image baseline}
\label{sec:baseline}

The baseline follows the project brief: the complete field photograph, without frame cropping,
is resized to $224\times224$ and passed through a frozen DINOv3 ViT-S/16~\cite{dinov3}. Its 384-dimensional \texttt{[CLS]}
embedding goes to a regression head
$\text{Linear}(384,256)\!\to\!\text{ReLU}\!\to\!\text{Dropout}(0.3)\!\to\!\text{Linear}(256,1)$
with a sigmoid scaled to $[0,100]$\,\%. At this resolution a seedling covers only a few pixels,
which motivates the plant-level approach below. The baseline is trained with the Huber loss
and a single seed (42).

\subsection{Multiple instance learning with area-weighted pooling}
\label{sec:mil}

We treat an image as a bag of plant regions (instances) with a single bag label $y$, the expert
score. Plants have no labels of their own, and we do not copy $y$ onto each plant. Every region is
resized to $224\times224$ and embedded with the same frozen DINOv3 backbone, giving features
$h_i \in \mathbb{R}^{384}$. Since the backbone is frozen, all bags are embedded once and cached.
The head $f$ from Section~\ref{sec:baseline} predicts a damage score $\hat y_i = f(h_i)$ for each
plant, and the image prediction is the visible-area-weighted mean
\begin{equation}
  \hat y = \sum_i w_i\, \hat y_i, \qquad w_i = \frac{A_i}{\sum_j A_j},
  \label{eq:weighted}
\end{equation}
as suggested by the supervisor, since the expert estimates the damaged share of the visible leaf
area. Large plants therefore count more, and small false-positive regions have little influence.

We compare this rule with three alternatives. \emph{Uniform} pooling sets $w_i = 1/N$.
\emph{ABMIL}~\cite{ilse2018attention} learns the weights from the features,
$a_i = \operatorname{softmax}_i\!\big(w^\top \tanh(V h_i)\big)$, and regresses from the pooled
feature $\sum_i a_i h_i$. \emph{Gated ABMIL} multiplies the $\tanh$ branch by a sigmoid gate
$\sigma(U h_i)$. The attention width is 128.

All MIL models are trained with the Huber loss, AdamW (learning rate $10^{-3}$, weight decay
$10^{-4}$), batch size 32, and at most 50 epochs with early stopping on validation MAE (patience
10). Each aggregation method is trained with seeds 42, 43 and 44.

\paragraph{Ranking-based training.} The expert scores are subjective, but the \emph{order} of two
images is often clear even when the exact values are not. Following the ranking direction of the
project brief, we train the area-weighted model on pairs $(A,B)$: every training image $A$ is paired
with a random partner $B$ whose score differs by at least $m$, and the loss is
\begin{equation}
  \mathcal{L} = \tfrac12\big[\ell_\delta(\hat y_A, y_A) + \ell_\delta(\hat y_B, y_B)\big]
  + \lambda \max\!\big(0,\, -s\,(\hat y_A - \hat y_B) + m\big),
  \qquad s = \operatorname{sign}(y_A - y_B),
\end{equation}
where $\ell_\delta$ is the Huber loss ($\delta = 1$), $m = 5$ percentage points and
$\lambda = 0.5$. Backbone, pooling and head are identical to the regression model, so only the
training objective differs.

\subsection{Out-of-distribution evaluation}
\label{sec:ood-method}

To test generalization, we apply all trained models zero-shot, without any fine-tuning or
recalibration, to three sets that were never used in development: (i) the field trial Weilburger Grenze
(JLU WG1, 07.10.2025, BBCH10, 906 scored images); (ii) the field trial DSV Asendorf (Trial~1,
15.--17.09.2025, BBCH11, 897 images); and (iii) the 218 Gro\ss-Gerau images from plots outside the
training and validation splits on which JLU and GAU disagreed by more than 5 points. The last set
shares field, date and camera with the training data but has noisy labels, and it lets us compare
model--rater agreement with rater--rater agreement. The pipeline is run unchanged, and we record
whether the frame was found and how many plant regions were extracted in each image.

\subsection{Metrics}
\label{sec:metrics}

We report MAE and RMSE in percentage points, Pearson $r$ and Spearman $\rho$. Because breeders
mainly need to \emph{rank} plots, we also report pairwise ranking accuracy: the fraction of image
pairs whose true scores differ by at least 5 points and whose predicted order is correct. For the
test set we give a bootstrap 95\,\% confidence interval of $\rho$ (5{,}000 resamples). Constant
mean and median predictors serve as reference points.
